\subsection{W-13 (SWI)}

Konstrukce, objektové paradigma, základní pravidla návrhu (SRP, LSP, DRY), refactoring (příznaky a vybrané techniky).

\subsubsection*{Objektové paradigma}

\begin{itemize}
  \item Systém je dekomponován na objekty, které spolu komunikují.
  \item Každý objekt má vlastní data a metody, které realizují zodpovědnost za vykonání určité funkcionality.
  \item Třída je šablona podle které jsou objekty vytvářeny.
\end{itemize}

Principy OOP
\begin{itemize}
  \item Zapouzdření - skrytí implementace objektu před okolím.
  \item Dědičnost - schopnost třídy dědit vlastnosti a metody od jiné třídy.
  \item Abstraktní třídy
  \item Polymorfismus - při volání metody je implementace určena typem objektu, ne typem proměnné.
\end{itemize}

Základní pravidla návrhu:
\begin{itemize}
  \item \textbf{Single Responsibility Principle (SRP)} - třída by měla mít pouze jednu odpovědnost.
  \item \textbf{Liskov Substitution Principle (LSP)} - objekty základní třídy by měly být zaměnitelné za objekty odvozené třídy.
  \item \textbf{Don't Repeat Yourself (DRY)} - opakující se kód by měl být extrahován do samostatné metody nebo třídy.
\end{itemize}
